\subsection{Change of base rings.}

In this section, we denote by A, B, A', B' four rings, by h and h' homomorphisms from A to A' and from B to B' respectively, by G an (A, B)-bimodule, by G' an (A', B')-bimodule, and by u a homomorphism from the underlying abelian group of G to the underlying abelian group of G', satisfying

\[
u (a g b) = h (a) u (g) h ^ {\prime} (b) \quad (a \in \mathrm{A}, g \in \mathrm{G}, b \in \mathrm{B}).\tag{11}
\]

Let E (resp. F) be a left A-module (resp. a right B-module). Recall (Chap. III, 2nd ed., App. II, no. 10) that, if one regards A′ (resp. B′) as a right A-module (resp. left B-module), the tensor product \(E' = A' \otimes_{A} E\) (resp. \(F' = F \otimes_{B} B'\) ) is endowed with a left A'-module (resp. right B'-module) structure defined by

\[
\begin{array}{c c} a _ {1} ^ {\prime} (a ^ {\prime} \otimes x) = (a _ {1} ^ {\prime} a ^ {\prime}) \otimes x & (a ^ {\prime}, a _ {1} ^ {\prime} \in A ^ {\prime}, x \in E) \\ (\text { resp. } (y \otimes b ^ {\prime}) b _ {1} ^ {\prime} = y \otimes (b ^ {\prime} b _ {1} ^ {\prime}) & (b ^ {\prime}, b _ {1} ^ {\prime} \in B ^ {\prime}, y \in F)). \end{array}\tag{12}
\]

PROPOSITION 1. --- Let E be a left A-module and F a right B-module; set \(\mathbf{E}^{\prime} = \mathbf{A}^{\prime}\otimes_{\mathbf{A}}\mathbf{E}\) and \(\mathbf{F}^{\prime} = \mathbf{F}\otimes_{\mathbf{B}}\mathbf{B}^{\prime}\) . For every bilinear map \(\Phi\) of \(\mathbf{E}\times \mathbf{F}\) into G, there exists one and only one bilinear map \(\Phi^{\prime}\) from \(\mathbf{E}^{\prime}\times \mathbf{F}^{\prime}\) into G' such that one has

\[
\Phi^ {\prime} (a ^ {\prime} \otimes x, y \otimes b ^ {\prime}) = a ^ {\prime}. u (\Phi (x, y)). b ^ {\prime}\tag{13}
\]

for all \(a' \in \mathrm{A}'\) , \(b' \in \mathrm{B}'\) , \(x \in \mathrm{E}\) , \(y \in \mathrm{F}\) .

The uniqueness of \(\Phi'\) follows from the fact that the elements \(a' \otimes x\) and \(y \otimes b'\) generate E' and F' respectively. To prove their existence, consider the map

\[
m: (a ^ {\prime}, x, y, b ^ {\prime}) \rightarrow a ^ {\prime}. u (\Phi (x, y)). b ^ {\prime}
\]

from A' × E × F × B' into G; it is evidently Z-multilinear, and it satisfies

\[
\begin{array}{r} m (a ^ {\prime}, a x, y, b ^ {\prime}) = m (a ^ {\prime} h (a), x, y, b ^ {\prime}) \\ m (a ^ {\prime}, x, y b, b ^ {\prime}) = m (a ^ {\prime}, x, y, h ^ {\prime} (b) b ^ {\prime}) \\ (a \in \mathrm{A}, b \in \mathrm{B}, a ^ {\prime} \in \mathrm{A} ^ {\prime}, b ^ {\prime} \in \mathrm{B} ^ {\prime}, x \in \mathrm{E}, y \in \mathrm{F}). \end{array}
\]

There therefore exists a Z-bilinear map \(\Phi'\) from E' × F' into G' satisfying (13) (Chap. III, 2nd ed., App. II, no. 1, prop. 2). This relation and the definition of the module structures of E' and F' by (12) show that \(\Phi'\) is bilinear, which completes the proof.

Given the hypotheses and notation of Proposition 1, let us now study the linear maps associated with \(\Phi\) and with \(\Phi'\) (no. 1, def. 2). To this end, we shall first define a canonical homomorphism from \(\mathfrak{L}_{A}(\mathrm{E}, \mathrm{G})\) into \(\mathfrak{L}_{A'}(\mathrm{E}', \mathrm{G}')\) . For every \(\nu \in \mathfrak{L}_{\Lambda}(\mathrm{E}, \mathrm{G})\) the map \((a', x) \to a'. u(\nu(x))\) of \(\mathrm{A}' \times \mathrm{E}\) into \(\mathrm{G}'\) is Z-bilinear, and, in view of (11), maps \((a'h(a), x)\) and \((a', ax)\) ( \(a \in \mathrm{A}\) ) onto the same element of \(\mathrm{G}'\) ; it therefore defines (Chap. III, 2nd ed.,

App. II, nos. 1 and 10) a map \(k(\nu)\) of \(\mathbf{E}' = \mathbf{A}' \otimes_{\mathbf{A}} \mathbf{E}\) into \(\mathbf{G}'\) such that \(k(\nu)(a' \otimes x) = a'. u(\nu(x))\) , and which, by (12), is A'-linear. Moreover, one deduces immediately from (12) that the map \(\nu \to k(\nu)\) of \(\mathfrak{L}_{\mathbf{A}}(\mathbf{E}, \mathbf{G})\) into \(\mathfrak{L}_{\mathbf{A}'}(\mathbf{E}', \mathbf{G}')\) satisfies \(k(\nu b) = k(\nu) h'(b)\) for all \(b \in B\) . Let \(i\) denote the canonical map \(y \to y \otimes 1\) of F into \(\mathbf{F}'\) . Then the diagram

\[
\begin{array}{c c c} \mathrm{F} & \xrightarrow {d _ {\Phi}} & \mathcal {L} _ {A} (\mathrm{E}, \mathrm{G}) \\ \downarrow i & & \downarrow k \\ \mathrm{F} ^ {\prime} & \xrightarrow {d _ {\Phi^ {\prime}}} & \mathcal {L} _ {A^{\prime}} (\mathrm{E} ^ {\prime}, \mathrm{G} ^ {\prime}) \end{array}\tag{14}
\]

(where \(d_{\Phi}\) and \(d_{\Phi'}\) denote the linear maps associated on the right to \(\Phi\) and \(\Phi'\) ) is commutative. Indeed, for \(x \in \mathrm{E}, y \in \mathrm{F}\) and \(a' \in \mathrm{A}'\) , one has \(d_{\Phi'}(i(y))(a' \otimes x) = \Phi'(a' \otimes x, y \otimes 1) = a'. u(\Phi(x, y)) = a'. u(d_{\Phi}(y)(x))\) , that is to say \(d_{\Phi'}(i(y))(a' \otimes x) = k(d_{\Phi}(y))(a' \otimes x)\) . One has an analogous commutation relation for the linear maps \(s_{\Phi}\) and \(s_{\Phi'}\) associated on the left to \(\Phi\) and \(\Phi'\) .

PROPOSITION 2. --- Suppose that B and B' are equipped with anti-automorphisms J and I such that

\[
h ^ {\prime} (b ^ {J}) = h ^ {\prime} (b) ^ {I} \quad \text { for all } b \in \mathrm{B}.\tag{15}
\]

Let E be a left A-module and F a left B-module; set \(E' = A' \otimes_{A} E\) and \(F' = B' \otimes_{B} F\) . For every sesquilinear map (with respect to J) \(\Phi\) of \(E \times F\) into G, there exists one and only one sesquilinear map (for I) \(\Phi'\) from \(E' \times F'\) into \(G'\) such that

\[
\Phi^ {\prime} (a ^ {\prime} \otimes x, b ^ {\prime} \otimes y) = a ^ {\prime}. u (\Phi (x, y)). b ^ {\prime I}\tag{16}
\]

for all \(a' \in \mathbf{A}'\) , \(b' \in \mathbf{B}'\) , \(x \in \mathbf{E}\) , \(y \in \mathbf{F}\) .

The uniqueness of \(\Phi'\) follows from the fact that the tensor products \(a' \otimes x\) and \(b' \otimes y\) generate E′ and F′ respectively. To establish their existence, consider the map

\[
m: (a, x, b ^ {\prime}, y) \rightarrow a ^ {\prime}. u (\Phi (x, y)). b ^ {\prime I}
\]

of \(\mathbf{A}' \times \mathbf{E} \times \mathbf{B}' \times \mathbf{F}\) into \(\mathbf{G}\). It is evidently \(\mathbf{Z}\)-multilinear, and, taking (11) and (15) into account, satisfies \(m(a', ax, b', y) = m(a'h(a), x, b', y)\) and \(m(a', x, b', by) = a'. u(\Phi(x, y)) \cdot h'(b)^{I} b'^{I} = m(a', x, b'h'(b), y)\) ( \(a \in \mathbf{A}, b \in \mathbf{B}, a' \in \mathbf{A}'\) , \(b' \in \mathbf{B}'\) , \(x \in \mathbf{E}\) , \(y \in \mathbf{F}\) ). There therefore exists a \(\mathbf{Z}\)-bilinear map \(\Phi'\) from \(\mathbf{E}' \times \mathbf{F}'\) into \(\mathbf{G}'\)

satisfying (16) (chap. III, 2nd ed., App. II, no. 1, prop. 2). This relation, together with the definition of the module structures of E′ and F′ by (12), shows, in view of (15), that Φ′ is sesquilinear for I, which completes the proof.

The most important examples of \((\mathbf{A}^{\prime}, \mathbf{B}^{\prime})\) -bimodules \(G^{\prime}\) equipped with Z-linear maps u from G into \(G^{\prime}\) satisfying (11), are the following:

\begin{enumerate}
\def\labelenumi{\arabic{enumi})}
\tightlist
\item
  We take for G′ the tensor product A′ ⊗\_A G ⊗\_B B′ (chap. III, 2nd ed., App. II, no. 9) and for u the map
\end{enumerate}

\[
g \to 1 \otimes g \otimes 1\tag{\(g\in G\)}
\]

from G to G'. The pair (G', u) thus defined is visibly universal in the following sense: for every (A', B')-bimodule G₁' and every Z-linear map u₁ from G to G₁' satisfying the analogue of (11), there exists a unique Z-linear map f from G' to G₁' such that f(a'g'b') = a'f(g')b' (a' ∈ A', g' ∈ G', b' ∈ B'; in other words, f is a homomorphism for the bimodule structures of G' and G₁') and such that u₁ = f ∘ u.

\begin{enumerate}
\def\labelenumi{\arabic{enumi})}
\setcounter{enumi}{1}
\item
  When A = B = G (the structure of (A, A)-bimodule of A being defined by left and right homotheties), \(A' = B'\) , and \(h = h'\) one can take for \(G'\) the ring \(A'\) and for u the homomorphism h from A into \(A'\) .
\item
  Suppose that we have A = B, \(A' = B'\) , \(h = h'\) , that the rings A and \(A'\) are commutative, and that the structure of left A-module of G coincides with its structure of right A-module. One can then take for \(G'\) the tensor product \(A' \otimes_{A} G\) (the structure of \(A'\) -module on the right of \(G'\) coinciding with its structure of \(A'\) -module on the left) and for u the map \(g \to 1 \otimes g\) ( \(g \in G\) ) from G into \(G'\) . We shall then say that the bilinear (resp. sesquilinear) map \(\Phi'\) defined by prop. 1 (resp. prop. 2) is obtained from \(\Phi\) by extension of the base ring, or by extension of scalars.
\end{enumerate}

What follows is valid equally for bilinear maps and for sesquilinear maps; the hypotheses and notations are those of prop. 1 (resp. prop. 2). Given a submodule M of E or F, we denote by M' the submodule of E' or F' generated by the canonical image of M.

PROPOSITION 3. --- With the hypotheses and notations of prop. 1 (resp. prop. 2), suppose moreover that A, B, A', B' are fields and that the maps \(\alpha\) and \(\beta\) from A'⊗A G and G⊗B B' into G' characterized by \(\alpha(a'\otimes g)=a'u(g)\) and \(\beta(g\otimes b')=u(g)b'\) ( \(a'\in A'\) , \(b'\in B'\) , \(g\in G\) ) are injective. Let M be a subspace of E and N a subspace of F. Then the subspace \((M')^{0}\) of F' orthogonal to M' with respect to \(\Phi'\) is equal to \((M^{0})^{\prime}\) , and, similarly, we have \((N')^{0} = (N^{0})^{\prime}\) .

Indeed, the inclusions \((\mathbf{M}^{0})^{\prime}\subset(\mathbf{M}^{\prime})^{0}\) and \((\mathbf{N}^{0})^{\prime}\subset(\mathbf{N}^{\prime})^{0}\) are evident (and moreover true without hypotheses on A, B, A', B', \(\alpha\) or \(\beta\) ). We shall prove the inclusion \((\mathbf{M}^{\prime})^{0}\subset(\mathbf{M}^{0})^{\prime}\) ; we leave to the reader the verification of the inclusion \((\mathbf{N}^{\prime})^{0}\subset(\mathbf{N}^{0})^{\prime}\) , which is entirely analogous. Let \(y^{\prime}\) be an element of \((\mathbf{M}^{\prime})^{0}\) . One can write

\[
y ^ {\prime} = \sum_ {i = 1} ^ {s} y _ {i} \otimes b _ {i} ^ {\prime} \quad (\text { resp. } y ^ {\prime} = \sum_ {i = 1} ^ {s} b _ {i} ^ {\prime} \otimes y _ {i})
\]

where \(y_i \in \mathbf{F} (1 \leqslant i \leqslant s)\) , and where the \(b_i'\) are elements of \(\mathbf{B}'\) which are linearly independent over \(\mathbf{B}\) for the structure of \(\mathbf{B}\) -module on the left (resp. on the right) of \(\mathbf{B}'\) . Let \(x \in \mathbf{M}\) and \(x' = 1 \otimes x \in \mathbf{M}'\) . We have

\[
0 = \Phi^ {\prime} (x ^ {\prime}, y ^ {\prime}) = \sum_ {i} u (\Phi (x, y _ {i})) b _ {i} ^ {\prime} = \beta (\sum_ {i} \Phi (x, y _ {i}) \otimes b _ {i} ^ {\prime})
\]

\[
(\text { resp. } 0 = \Phi^ {\prime} (x ^ {\prime}, y ^ {\prime}) = \sum_ {i} u (\Phi (x, y _ {i})) b _ {i} ^ {\prime \mathrm{I}} = \beta (\sum_ {i} \Phi (x, y _ {i}) \otimes b _ {i} ^ {\prime \mathrm{I}}).
\]

Since β is injective and since the \(b_i'\) (resp. the \(b _ {i} ^ {\prime \mathrm{I}}\) , taking (15) into account) are linearly independent over B for the structure of left B-module of B', this entails \(\Phi(x, y_i) = 0\) for \(i = 1, \ldots, s\) . Since this last relation holds for every \(x \in \mathbf{M}\) , one has \(y_i \in \mathbf{M}^0\) for \(i = 1, \ldots, s\) , whence \(y' \in (\mathbf{M}^0)'\) . QED.

COROLLARY. --- With the hypotheses and notations of prop. 3, for \(\Phi'\) to be nondegenerate, it is necessary and sufficient that \(\Phi\) be nondegenerate.

Indeed, by prop. 3, we have \((\mathbf{F}')^0 = (\mathbf{F}^0)'\) and \((\mathbf{E}')^0 = (\mathbf{E}^0)'\) . On the other hand, for \(\Phi\) (resp. \(\Phi'\) ) to be nondegenerate, it is necessary and sufficient that one have \(\mathbf{F}^0 = \mathbf{E}^0 = \{0\}\) (resp. \((\mathbf{F}')^0 = (\mathbf{E}')^0 = \{0\}\) ).

Remark. --- Suppose that A, B, A' and B' are fields. Then, for the bimodules G' defined in the three examples above, the maps α and β are injective, as follows immediately from chap. III, 2nd ed., App. II, no. 6.

